\section{Critical Pairs as Generators of Holonomy}
\label{sec:critical-pairs-curvature}

In string and term rewriting, the one-step relation is generated by finitely many rules applied anywhere, so a single word or term typically has many peaks. Elementary flatness asks for all of them to be joinable. This section recalls why only a small, structured part of them can fail, namely the overlaps recorded by critical pairs, and explains how we use critical pairs relative to chosen start terms.

\subsection{Three kinds of peaks}
\label{subsec:three-kinds-of-peaks}

Consider a term $t$ with two redexes, one for a rule $l_1 \to r_1$ at position $p$ and one for a rule $l_2 \to r_2$ at position $q$. There are three possibilities \cite{Huet1980,BaaderNipkow1998,Terese2003}, illustrated in \Cref{fig:three-peaks}.

\begin{enumerate}
    \item \emph{Disjoint redexes.} Neither position lies below the other. The two rewrites act on separate parts of $t$ and commute: each result can be rewritten by the other rule to the same term. For example, with rules $a \to b$ and $c \to d$, the word $ac$ has the peak $bc \leftarrow ac \to ad$, joined at $bd$. No two left-hand sides overlap here.
    \item \emph{Variable overlap.} One redex lies inside the part of the other that is matched by a variable. The peak is again joinable: on the side where the outer rule has fired, the inner rewrite is applied to every copy of that variable's instance on the right-hand side; on the other side, it is applied to any remaining copies on the left-hand side, after which the outer rule fires. For left-linear rules there are no remaining copies. For example, with rules $f(x) \to g(x)$ and $a \to b$, the term $f(a)$ has the peak $g(a) \leftarrow f(a) \to f(b)$, joined at $g(b)$.
    \item \emph{Critical overlap.} One redex lies at a non-variable position of the other's left-hand side, so the two left-hand sides genuinely share symbols. Joinability is not automatic. Every such peak is an instance, inside some context and under some substitution, of a \emph{critical pair}: a pair of terms obtained by unifying one left-hand side with a non-variable subterm of another, using a most general unifier. A rule may overlap with a renamed copy of itself at a proper subterm; the overlap of a rule with itself at the root is excluded because it produces two identical terms.
\end{enumerate}

\begin{figure}[t]
\centering
\begin{tikzpicture}
  \begin{scope}
    \node[st] (s) at (0,0) {$ac$};
    \node[st] (l) at (-0.9,-1.0) {$bc$};
    \node[st] (r) at (0.9,-1.0) {$ad$};
    \node[st] (w) at (0,-2.0) {$bd$};
    \begin{scope}[on background layer]\fill[cellfill] (s.center)--(l.center)--(w.center)--(r.center)--cycle;\end{scope}
    \draw[rw] (s)--(l); \draw[rw] (s)--(r); \draw[rw] (l)--(w); \draw[rw] (r)--(w);
    \node[font=\small] at (0,-2.75) {(a) disjoint};
    \node[font=\scriptsize] at (0,-3.25) {$a \to b$,\; $c \to d$};
  \end{scope}
  \begin{scope}[xshift=4.4cm]
    \node[st] (s) at (0,0) {$f(a)$};
    \node[st] (l) at (-0.9,-1.0) {$g(a)$};
    \node[st] (r) at (0.9,-1.0) {$f(b)$};
    \node[st] (w) at (0,-2.0) {$g(b)$};
    \begin{scope}[on background layer]\fill[cellfill] (s.center)--(l.center)--(w.center)--(r.center)--cycle;\end{scope}
    \draw[rw] (s)--(l); \draw[rw] (s)--(r); \draw[rw] (l)--(w); \draw[rw] (r)--(w);
    \node[font=\small] at (0,-2.75) {(b) variable overlap};
    \node[font=\scriptsize] at (0,-3.25) {$f(x) \to g(x)$,\; $a \to b$};
  \end{scope}
  \begin{scope}[xshift=8.8cm]
    \node[st] (s) at (0,0) {$aba$};
    \node[nf] (l) at (-0.9,-1.0) {$aa$};
    \node[st] (r) at (0.9,-1.0) {$ab$};
    \node[nf] (w) at (0.9,-2.0) {$a$};
    \draw[rw] (s)--(l); \draw[rw] (s)--(r); \draw[rw] (r)--(w);
    \draw[defect] (l) .. controls (-0.7,-1.9) and (0.1,-2.2) .. (w);
    \node[font=\small] at (0,-2.75) {(c) critical overlap};
    \node[font=\scriptsize] at (0,-3.25) {$ab \to a$,\; $ba \to b$};
  \end{scope}
\end{tikzpicture}
\caption{The three kinds of peaks. (a) Rewrites at disjoint positions commute. (b) A rewrite inside the instance of a variable is absorbed after the outer rule fires. Both are always joinable. (c) Overlapping left-hand sides: in $aba$, the redexes $ab$ and $ba$ share the middle letter. This is the string example $E006$; its branches lead to the distinct normal forms $aa$ and $a$ (boxed), so the branching stays open.}
\label{fig:three-peaks}
\label{fig:string-overlap-nonconfluent}
\end{figure}

String rewriting is the special case of term rewriting with unary function symbols, and the same classification applies to words. Two left-hand sides overlap critically when a suffix of one is a prefix of the other, or when one occurs inside the other.

The first two cases are always joinable, so a nonjoinable peak must come from a critical overlap. Conversely, if every critical pair is joinable, then so is every instance of it in any context, because rewriting is closed under contexts and substitutions. This is the content of the classical critical-pair lemma.

\begin{theorem}[Critical-pair lemma \cite{KnuthBendix1970,Huet1980}]
\label{thm:critical-pair-lemma}
A term rewriting system is locally confluent if and only if all of its critical pairs are joinable. The same holds for string rewriting systems.
\end{theorem}

We use this theorem as stated in the literature (see also \cite[Chapter~6]{BaaderNipkow1998}); it is neither re-proved nor mechanized here. Combined with \Cref{thm:elementary-flatness-iff-local-confluence} and \Cref{cor:terminating-confluence-iff-flat}, it gives the holonomy form of the Knuth--Bendix criterion.

\begin{corollary}
\label{cor:cp-flatness}
A term or string rewriting system is elementarily flat if and only if all of its critical pairs are joinable. If the system is terminating, it is confluent if and only if all of its critical pairs are joinable.
\end{corollary}

In the language of \Cref{sec:reduction-two-complex}, critical pairs generate the nontrivial holonomy obligations: they are the branchings whose fillers must be checked, while disjoint and variable-overlap peaks are filled automatically. They are not the whole set of peaks. Even a joinable critical pair is a genuine generating branching, whose $2$-cell happens to be filled.

\subsection{Critical pairs relative to start terms}
\label{subsec:string-critical-pairs}
\label{subsec:reachable-critical-pairs-vs-peaks}

Our computations study a rewriting system from designated start terms. Let $S$ be the set of starts and $X_S$ the set of terms reachable from $S$. Every descendant of a term in $X_S$ lies in $X_S$, so the rewriting relation restricts to $X_S$, and two terms of $X_S$ are joinable in the restricted system exactly when they are joinable in the full one.

\begin{definition}[Reachable critical-pair instance]
\label{def:reachable-cp-instance}
A \emph{reachable critical-pair instance} is a peak $u \leftarrow t \to v$ with $t \in X_S$ and $u \neq v$ in which the two redexes form a critical overlap. It is recorded by its source $t$, its unordered pair of branches $\{u,v\}$, and the rules and positions involved.
\end{definition}

Two features of this definition matter. First, the instance is concrete. The source of a critical pair as computed from the rules is a most general overlap, such as $ab$ for the rules $ab \to x$ and $ab \to y$, or $f(g(x))$ for $f(g(x)) \to p(x)$ and $g(y) \to q(y)$. The terms actually reached from the starts contain instances of it: the word $zab$, or the ground term $f(g(a))$. Asking whether the general source itself is reachable misses these. Second, the definition depends on $S$. A critical pair may have no instance in $X_S$, in which case it plays no role in the restricted system.

The classification of \Cref{subsec:three-kinds-of-peaks} applies peak by peak, and the joins of disjoint and variable-overlap peaks stay inside $X_S$. Hence the restricted system on $X_S$ is locally confluent if and only if every reachable critical-pair instance is joinable. This is the start-relative form of \Cref{thm:critical-pair-lemma} that the computations test. The equivalence holds whether or not $X_S$ is finite. Our computational certification of it, however, requires complete exploration of $X_S$: the absence of a join witness in an incomplete bounded search does not certify nonjoinability (\Cref{subsec:certification}).

\subsection{Curated examples}
\label{subsec:left-linear-trs-bridge}

\Cref{tab:system-vs-reachable-cp-vs-graph-peak} compares, for the curated string and term examples, the number of critical pairs computed from the rules, the number of reachable critical-pair instances, and the number of peaks in the explored reduction graph. Since disjoint and variable-overlap peaks are joinable, the general relation is only that every \emph{nonjoinable} peak is a critical-pair instance. In the curated examples something stronger happens: every peak of the explored graph is a critical-pair instance and conversely. This is a property of these examples, which happen to contain no disjoint or variable-overlap peaks, and not a general fact.

\begin{table}[t]
    \caption{Critical pairs computed from the rules (``system''), reachable critical-pair instances, peaks in the explored reduction graph, and nonjoinable reachable instances. In every example the explored graph is complete, and the reachable instances coincide with the peaks.}
    \label{tab:system-vs-reachable-cp-vs-graph-peak}
    \centering
    \smallskip
    \small
\setlength{\tabcolsep}{4pt}
\begin{tabularx}{\linewidth}{@{}L{0.22\linewidth} L{0.11\linewidth} *{4}{>{\centering\arraybackslash}p{0.08\linewidth}} >{\centering\arraybackslash}p{0.11\linewidth}@{}}
\toprule
Example & Domain & \shortstack{System\\CPs} & \shortstack{Reach.\\CPs} & \shortstack{Graph\\peaks} & \shortstack{Defective\\CPs} & Confluent? \\
\midrule
$E005$ (sort)          & string    & 1 & 1 & 1 & 0 & yes \\
$E006$ (overlap)       & string    & 2 & 1 & 1 & 1 & no \\
$E007$ (before)        & string    & 1 & 1 & 1 & 1 & no \\
$E008$ (after)         & string    & 1 & 1 & 1 & 0 & yes \\
$TRS002$ (root)        & term rew. & 1 & 1 & 1 & 1 & no \\
$TRS003$ (nested)      & term rew. & 1 & 1 & 1 & 1 & no \\
$TRS004$ (resolved)    & term rew. & 1 & 1 & 1 & 0 & yes \\
\bottomrule
\end{tabularx}

\end{table}

The string example $E006$ (\Cref{fig:three-peaks}c) has rules $ab \to a$ and $ba \to b$, and the rules produce two critical pairs, from the overlaps $aba$ and $bab$. From the start word $aba$ the reachable words are $aba$, $aa$, $ab$, and $a$. Only the first critical pair has an instance there, and it is the single peak of the explored graph. Its branches $aa$ and $ab \to a$ end in different normal forms, so it is nonjoinable and the system is not confluent from $aba$. The second critical pair is irrelevant from this start.

\begin{observation}[Curated string examples]
\label{prop:reachable-critical-pairs-capture-string-obstructions}
In each of the string examples $E005$--$E008$, explored completely from its start word, the reachable critical-pair instances are exactly the peaks of the explored reduction graph, the nonjoinable instances are exactly the nonjoinable peaks, and the system is confluent from its start word exactly when there are none.
\end{observation}

The supporting computations are in \Cref{tab:system-vs-reachable-cp-vs-graph-peak,tab:curated-string-summary}. Each example is small enough to check by hand. In $E005$, the sorting rules $ba \to ab$, $ca \to ac$, $cb \to bc$ give, from $cba$, a single overlap of $cb$ and $ba$; both branches $bca$ and $cab$ reach $abc$. In $E006$, as above, the single reachable instance is nonjoinable. The pair $E007$--$E008$ is treated in \Cref{sec:completion-holonomy}.

The term examples test the same picture with first-order terms. In $TRS002$, the rules $f(a,y) \to g(y)$ and $f(x,b) \to h(x)$ overlap at the root of $f(a,b)$, giving the branches $g(b)$ and $h(a)$, which are distinct normal forms. In $TRS003$, the rules
\[
f(g(x)) \to p(x), \qquad g(a) \to b
\]
overlap below the root: the redex $g(a)$ sits at the argument position of $f(g(x))$, a non-variable position. The reachable instance has source $f(g(a))$ and branches $p(a)$ and $f(b)$, again distinct normal forms. In $TRS004$ the rule $f(b) \to p(a)$ is added; the same instance remains, but now $f(b) \to p(a)$ fills it.

\begin{observation}[Curated left-linear term examples]
\label{prop:trs-critical-pairs-as-curvature-generators}
In each of the term examples $TRS002$--$TRS004$, explored completely from its start term, the reachable critical-pair instances are exactly the peaks of the explored reduction graph, and the system is confluent from its start term exactly when every instance is joinable. In $TRS003$ and $TRS004$ the overlap lies below the root.
\end{observation}

These observations are small. Their purpose is to confirm on concrete inputs that the implementation finds the right objects: instances in context rather than general overlap sources, overlaps below the root as well as at it, and a verdict tied to joinability.
